\section{Beyond a distance-only fill law}
\label{sec:realisticquoting}

A dealer HJB is an optimization equation, not a unique specification of
market behavior. The omitted mechanisms belong to a particular state,
event law and objective. The original Avellaneda--Stoikov model uses
exponential utility, not merely a quadratic inventory penalty
\citep{AvellanedaStoikov2008}. A formula involving \(1/k\) comes from a
particular exponential intensity cell; it is not a universal competitive
spread. For example, maximizing
\(Ae^{-k\delta}(\delta-c)\) gives the unconstrained interior condition
\(\delta=c+1/k\) only when the continuation-adjusted charge \(c\) is held
fixed and the quote is continuous. Tick constraints, endogenous marks,
queue resets and competitors can change that calculation.

\subsection{The six requirements for a continuing quote}
The modelling requirements of \citet[Section~37.1]{NoguerAlonso2026Execution} take a
specific form here. \emph{Fill law}: tagged rank, partial fills, reservations
and pending cancellation are explicit in Sections~3.3, 3.5 and~3.7; venue
priority rules and the joint law of feed and order latency are not.
\emph{Information}: the marked-wealth charge \eqref{eq:jointmarkedcharge} must use the joint law of fill
and price mark, and Proposition~\ref{prop:filljumpfrechet} shows exactly what
is lost when only its marginals are known. \emph{Objective}: the certificates
concern the additive marked-wealth objective; drawdown, VaR or Sharpe criteria
need a wealth-and-peak state or a separate formulation, and an adaptive
mean--variance criterion needs a precommitment or equilibrium interpretation.
Impact and liquidity, competition and evaluation carry over as stated there; a
reduced-form response to rival depth is not a Nash equilibrium, and a replay
that holds the historical book fixed cannot validate the book's response to a
different quoting policy.

\subsection{One event, one marked-wealth charge}

Let a joint event mark contain executed quantity \(\xi\) (positive for a
buy), execution price \(P\), the reference-price jump \(J\), cash fee
\(f\) (negative for a rebate), and any queue or message update. With
pre-event inventory \(q\) and reference \(S\), the exact wealth increment is
\begin{equation}\label{eq:jointmarkedcharge}
 \Delta W=\xi(S-P)+(q+\xi)J-f.
\end{equation}
Indeed cash changes by \(-\xi P-f\), and marked inventory changes from
\(qS\) to \((q+\xi)(S+J)\). Thus a finite-activity event kernel
\(\nu_t(\dd z\mid x,a)\), a between-event reward rate \(r_0\), and a
between-event generator \(\mathcal L_0\) give the schematic HJB
\begin{equation}\label{eq:jointmarkedhjb}
0=\partial_t v+\sup_{a\in\mathcal A_t(x)}
 \left\{r_0+\mathcal L_0v+
 \int[\Delta W(x,a,z)+v(t,F(x,a,z))-v(t,x)]
          \nu_t(\dd z\mid x,a)\right\}.
\end{equation}
Here \(v\) values \emph{future increments}, excluding already booked
wealth; running penalties belong in \(r_0\). This convention prevents
subtracting a fill markout twice. If a reward instead books a price mark
at a future horizon, the overlapping continuation must be adjusted, or
the state must retain the outstanding markout obligation. The shared
Bellman step requires a consistent accounting boundary.

\subsection{Sharp ambiguity from missing fill--jump dependence}

Consider one ask quote at spread \(d\), initial inventory zero, no
continuation or terminal inventory penalty, and a fee \(f\) per fill.
Let \(X\in\{0,1\}\) indicate its fill and \(J=hY\), where
\(Y\in\{0,1\}\), \(h>0\), is a simultaneous upward price jump. A
nonparticipating dealer earns zero in this cell. Its quoting gain is
\[
 G=\E[X(d-f-hY)].
\]
The endpoint liquidation convention is marked wealth, not cost-free
future liquidation. It is deliberately the same for both actions.

\begin{proposition}[A sharp marginal-only quoting interval]
\label{prop:filljumpfrechet}
If \(\Pr(X=1)=p\) and \(\Pr(Y=1)=q\), but their dependence is unknown,
then the exact set of possible quoting gains is
\begin{equation}\label{eq:frechetquoting}
 \left[p(d-f)-h\min(p,q),\quad
 p(d-f)-h\max(0,p+q-1)\right].
\end{equation}
Independence selects the single value \(p(d-f)-hpq\). A nonnegative
lower endpoint is necessary and sufficient for this quote to weakly
improve on nonparticipation for every joint law with these marginals.
\end{proposition}
\begin{proof}
Write \(z=\Pr(X=1,Y=1)\). The four cell probabilities are
\(z,p-z,q-z,1-p-q+z\). Their nonnegativity is equivalent to
\(\max(0,p+q-1)\le z\le\min(p,q)\). Every point in this interval is a
joint law, and gain is the decreasing affine function
\(p(d-f)-hz\). Its endpoint values give the interval and the comparison
with zero.
\end{proof}

The argument is the two-point marginal-coupling calculation of
\citet[Proposition~37.1]{NoguerAlonso2026Execution}; the proposition records its
market-making instance, in which the participation decision rather than a
routing decision is reversed.

Take \(p=0.40\), \(q=0.25\), \(d=1\), \(f=0\), \(h=2\), in tick
units. Independent marginals suggest a gain of \(0.20\) ticks per
opportunity. The admissible joint law with all upward jumps occurring
at fills has gain \(-0.10\). Fill probability and unconditional jump
probability are identical. Their full ambiguity interval is
\([-0.10,0.40]\). With probability \(0.10+0.15\theta\) assigned to the
joint fill/jump event, gain is \(0.20-0.30\theta\): the quote becomes
unprofitable once \(\theta>2/3\).

\begin{table}[htbp]
\centering\small
\caption{Dependence changes participation despite fixed marginal fill and
jump probabilities. Values are synthetic ticks per opportunity.}
\begin{tabular}{lrr}
\toprule
Law & Quoting gain & Selected action\\
\midrule
\inputtable{tables/realism_rows.tex}
\bottomrule
\end{tabular}
\end{table}

For partial fills and several marks, a finite coupling matrix gives the
same kind of exact stage calculation. Given marginal weights \(p_i,q_j\)
and a continuation-aware event gain \(g_{ij}\), its worst expectation is
\begin{equation}\label{eq:marktransport}
 \min_{\pi\ge0}\sum_{ij}\pi_{ij}g_{ij},\qquad
 \sum_j\pi_{ij}=p_i,\quad \sum_i\pi_{ij}=q_j.
\end{equation}
The event destinations and continuation enclosure must be included in
\(g_{ij}\); optimizing only instantaneous spread is insufficient. This is
a finite marginal-coupling linear program, not a new solution of the
full dealer game. To feed a stage bound into
Section~\ref{sec:baselinequoting}, retain the same underlying model in
candidate and baseline advantages. Independent minimization of their
values can discard cancellation of common uncertainty.

The script \texttt{verification/realism\_stress.py} checks the interval
against independent linear programs over a probability grid, including
degenerate marginals, and verifies the mixture switch. These exact toy
laws demonstrate an identification obstruction; they are not inferred
from market observations.

\subsection{What would make the added realism measurable?}

The evaluation unit should be an order/message episode with observable
queue history, submitted size, executions, cancel acknowledgements and
price marks at preregistered horizons. Preserve causality using receipt
times available to the agent. Split by time and keep overlapping marks
within the same partition. Estimate partial-fill distributions and the
joint fill/mark law conditional on rank, size, latency and regime, with
uncertainty increasing where action support is sparse. Compare a frozen
baseline with the same net fee and terminal inventory accounting.
Report net marked reward, inventory and drawdown tails, stale-quote
exposure, fill calibration and certificate coverage. An ablation that
adds each variable separately should be accompanied by combined latency,
toxicity and competitor-response stress. Neither a positive simulated
spread nor a narrow in-model interval supplies empirical profitability.
